\documentclass[11pt, a4paper]{article}

\usepackage[T1]{fontenc}
\usepackage[utf8]{inputenc}
\usepackage{tgpagella}
\usepackage[spanish, es-noshorthands]{babel}
\usepackage{amsmath, amssymb, bm}
\usepackage{tabularx}
\usepackage{booktabs}
\usepackage[most]{tcolorbox}
\usepackage[margin=2.5cm]{geometry}
\usepackage{ragged2e}
\usepackage{fancyhdr}
\usepackage[american]{circuitikz}

\pagestyle{fancy}
\fancyhf{}
\setlength{\headheight}{16pt}
\lhead{\textbf{UCB Sede Tarija} | Ing. Mecatrónica}
\chead{}
\rhead{IMT-121 Circuitos Electrónicos I | Guía Lab B}
\lfoot{Prof: M.Sc. Ing. Bernardo Quiroga Turdera}
\rfoot{Página \thepage}
\renewcommand{\headrulewidth}{0.4pt}

\definecolor{ucbblue}{RGB}{0, 51, 102}

\begin{document}

\begin{tcolorbox}[colback=ucbblue!5!white, colframe=ucbblue, title=\textbf{Experiencia Integrada B — Guía de Laboratorio}]
    \centering \large \textbf{Máxima Transferencia, Efecto de Carga y Curva de Potencia}
\end{tcolorbox}

\section{Propósito de la Experiencia}
Utilizar el circuito caracterizado en la Experiencia A para verificar físicamente el Teorema de Máxima Transferencia de Potencia. Se debe trazar experimentalmente la curva $P_L(R_L)$ y comprobar empíricamente el efecto de carga (loading effect).

\section{Preparación y Predicción Analítica}
No se puede "adivinar" el máximo. Su equipo debe usar los valores obtenidos en el laboratorio anterior (Teoría o SPICE si la medición aún no está estable) para predecir dónde ocurrirá el fenómeno.

\begin{itemize}
    \item \textbf{Tensión equivalente de su circuito:} $V_{th} = \text{\rule{3cm}{0.4pt}} \text{V}$.
    \item \textbf{Resistencia interna de su circuito:} $R_{th} = \text{\rule{3cm}{0.4pt}} \text{k}\Omega$.
    \item \textbf{Predicción del Máximo:} La teoría exige que $R_{L,\text{max}} = R_{th}$, por tanto, esperamos el pico de potencia en $R_{L,\text{pred}} = \text{\rule{3cm}{0.4pt}} \text{k}\Omega$.
\end{itemize}

\section{Selección de Cargas ($R_L$)}
Seleccione de su kit un rango de resistores que permita observar la campana completa. Si su $R_{th}$ es aproximadamente $2.2\,\text{k}\Omega$, un conjunto útil es:
\begin{center}
    \textbf{470\,$\bm{\Omega}$ - 1.0\,k$\bm{\Omega}$ - 1.5\,k$\bm{\Omega}$ - 2.2\,k$\bm{\Omega}$ - 3.3\,k$\bm{\Omega}$ - 4.7\,k$\bm{\Omega}$ - 10\,k$\bm{\Omega}$}
\end{center}
\textit{Nota: Debe tener valores claramente por debajo, uno muy cercano, y valores claramente por encima de su $R_{th}$.}

\section{Simulación en LTSpice (Barrido Paramétrico)}
A diferencia de un análisis `.dc` (que barre voltajes/corrientes de fuentes), para barrer valores de una resistencia utilizaremos el comando de pasos parametrizados. 
\begin{enumerate}
    \item Cambie el valor de su resistencia de carga a la variable \texttt{\{RL\}}.
    \item Agregue la directiva: \texttt{.param RL 2.2k}
    \item Agregue la directiva de barrido discreto: \\
    \texttt{.step param RL list 470 1k 1.5k 2.2k 3.3k 4.7k 10k}
    \item Mantenga el análisis en \texttt{.op}
    \item Registre en su tabla los valores de $V_L$ e $I_L$ arrojados por SPICE para cada paso de resistencia, y calcule $P_L = V_L \cdot I_L$.
\end{enumerate}

\section{Procedimiento Físico (Banco)}
\begin{enumerate}
    \item \textbf{Medición de Corriente (Opcional/Supervisada):} Si el instructor autorizó la medición de corriente en serie, registre $I_L$ directamente. De lo contrario, opte por el \textbf{Método Indirecto}: mida $V_L$ con el multímetro y calcule $I_L = V_L / R_{L,\text{medido}}$.
    \item Conecte sucesivamente cada una de las resistencias seleccionadas en el puerto $a-b$ de su circuito físico y anote los resultados en la Hoja de Datos.
    \item Analice el Efecto de Carga (Loading Effect): Observe cómo $V_L$ va cambiando a medida que aumenta la $R_L$.
\end{enumerate}

\end{document}
